\chapter{Estimation, teams, incidents, and engineering management}

\section{Estimation under uncertainty}
An estimate is a prediction with assumptions, not a promise extracted from a
person. Separate three quantities: effort, elapsed time, and confidence. A task
may require two days of work but a week of elapsed time because it depends on an
approval or a test environment.

Use ranges when uncertainty is material. State the reference class, known work,
unknowns, dependencies, and what would change the estimate. Decompose until a
piece is small enough to observe progress, but do not confuse more detailed
fiction with more accuracy.

\section{Planning and sequencing}
Sequence work by risk and learning, not only by apparent feature value. Resolve
unknown integrations early. Build a vertical slice that exercises the real path
from input to persistence and observation. This uncovers hidden work sooner than
building isolated layers that cannot yet run.

A useful plan has:

\begin{itemize}
  \item outcome and non-goals;
  \item decision log and assumptions;
  \item milestones that produce observable evidence;
  \item dependencies and owners;
  \item quality and operational exit criteria;
  \item a change and rollback path.
\end{itemize}

\section{Team organization}
Team structure changes communication paths and ownership. Conway's law is a
useful observation about correspondence between communication structures and
system design, but it is not a deterministic theorem. Align ownership with
cohesive capabilities where possible. Make interfaces and escalation paths
explicit when ownership must cross teams.

Avoid “one person owns quality.” Every engineer affects quality, while specific
roles may coordinate standards, testing, security, or operations. A team is
responsible for the system it changes; a platform or specialist team should
reduce friction, not become a queue for every decision.

\section{Engineering management}
Management is system design applied to people, information, incentives, and
constraints. Good managers make priorities, decision rights, and trade-offs
clear; protect focus; develop capability; and ensure that uncomfortable evidence
can travel upward. Metrics should support learning rather than invite gaming.

Useful signals include lead time, change failure rate, recovery time, defect
escape, review wait, and user outcomes, but each can be misinterpreted. DORA-style
delivery measures are indicators, not a scoreboard for comparing unlike teams.
Combine quantitative signals with qualitative context.

\section{Incidents and learning}
The same skills used in operational incidents apply to project incidents:
surface facts, assign a coordinator, stabilize, communicate, and record decisions.
A missed migration, privacy concern, or repeated failed release deserves a causal
review. A postmortem action should improve the system, workflow, or decision
boundary; it should not merely ask people to remember harder.

\section{Decision rights}
For a recurring decision, name who recommends, who decides, who must be
consulted, and who is informed. Centralize decisions that need consistency or
rare expertise; delegate reversible local decisions. Ambiguous authority creates
delay and makes risk invisible.

\section{Healthy delivery pressure}
Urgency is real, but speed without a definition of done is borrowing from the
future. When a deadline forces a shortcut, record what is deferred, what risk it
creates, and the condition for repayment. A safe temporary reduction in scope is
usually easier to recover from than an invisible reduction in reliability or
security.

\begin{exercise}
Create a one-page plan for a four-person team delivering a new order export.
Include outcome, non-goals, risks, milestones, quality gates, owner boundaries,
and a rollback plan. Include one decision that should be delegated and one that
should be escalated.
\end{exercise}

\section{Chapter summary}
Estimate with ranges and assumptions. Sequence for learning and risk reduction,
align ownership with capabilities, use metrics as evidence rather than targets,
and turn incidents and shortcuts into explicit system improvements.

